\section{Appendix: Scope and Qualifications}
\label{sec:scope}

{\small
\textbf{Edition, formulary and occurrence.} The study treats the Mass
\latin{Dominica decima octava post Pentecosten}, II classis, of the
\work{Missale Romanum}, editio typica, Typis Polyglottis Vaticanis 1962,
pp.~410--411, nos.~1669--1678, in the universal 1962 calendar and no particular
calendar. Its occurrence on 27~September 2026 rests on the Missal's own calendar
and general rubrics (RG 111~b, 127~b; RGMR 427, 434~b, 494~b); the dated 2026
Ordos of the Fraternity of St Peter in France and of the Institute of Christ the
King Sovereign Priest in France agree on the Sunday, its class and its colour,
and are cited for the date only. The particular solemnity of St Thérèse of the
Child Jesus that each prints belongs to those institutes and is not part of the
universal formulary. The precedence table of the general rubrics and the
Breviary rubric behind the Gloria, which the typical edition does not print,
were not reread.

\textbf{Text control.} Every Latin form of the ten proper elements was read on
the page images of the CMAA facsimile of the typical edition and compared with
the Benziger \latin{editio iuxta typicam} of 1962 (whose pointing and one
accent differ and are not followed) and with the Pustet Missal of 1862, in
which every element's wording already stands. The marginal numbers 1671--1674
are clipped in the facsimile and are fixed by continuity. The collation, the
conclusion forms and the comparison with the Clementine Vulgate are recorded in
\texttt{propers/verified.md}. The Latin source of the Introit's and the
Offertory's distinctive wording, and the psalter behind the Communion's
\latin{in aula sancta}, were not identified: the Vulgate and the Greek of the Septuagint account
for none of them, and no Roman Psalter or chant manuscript was consulted.

\textbf{English.} The scriptural English is the Douay--Rheims in Challoner's
revision at the canonical verses; where the Missal adapts its source (the
Introit's antiphon, the Epistle's address, the Alleluia's opening, the Gospel's
opening, the Communion's \latin{aula}, and the whole Offertory) that English
does not render the Missal's Latin, and no English has been supplied in its
place. The English of the three orations is the anonymous translation printed
by Cummiskey in Philadelphia in 1861, read in the library's tracked text of that
book and not on its page images. Neither English is an approved liturgical
translation of the 1962 Missal. The Trinity Preface is cited and not
reproduced. English renderings of patristic and medieval Latin that are not
drawn from a named translation are glosses of the Latin printed beside them.

\textbf{Reception.} Each appointed passage was searched in the Latin and Greek
Fathers and the later Doctors, saints and exegetes available to this
production, and the witnesses read are those named in the References at their
exact loci. The Greek of Chrysostom on Ps~121, Theodoret, the expositions
printed under St Athanasius's name and Theophylact was read in Migne's columns.
Chrysostom on Matthew, and Augustine on the psalms and the harmony of the
Gospels, are quoted from the CCEL text of the nineteenth-century Nicene and
Post-Nicene English, with Augustine's Latin read in Migne at the loci quoted;
Chrysostom on First Corinthians and Hebrews is quoted from New Advent's delivery
of the same translation, which modernises its second person.
Bellarmine was read only in O'Sullivan's abridged English of 1866, in a tracked
web transcription. Cornelius a Lapide was read in optical text layers of the
Antwerp printings, not on their page images, and is summarised rather than
quoted. St Hilary's commentary on Matthew was read on the page images of the
Maurist text as Migne reprints it; its modern critical edition was not reached.
St Ambrose and St Bede on Luke's telling of the healing were read in
transcriptions of Migne, not on page images. Aquinas's \work{Catena aurea} on
Matthew, which splices and rewords several authors at this Gospel, served only
to find leads, and each author is cited from his own work. No patristic commentary on Ecclus 36:18 or at the Offertory's cited
verses of Exodus was located; the Introit and the Offertory are read by the
liturgical commentators on the Mass. Anthony of Padua's
Sunday sermon on this Gospel was read in a modern translation and is summarised,
not quoted. The attribution of the continuation of \work{The Liturgical Year}
to Dom Lucien Fromage rests on the library's standing registry and the
initials that sign the French preface. No witness is cited for a reading of this
Mass as a whole except the liturgical commentators. Rupert of Deutz, Honorius,
Sicard, Durandus and Berno expounded books in which this Mass had another
Gospel, and they are cited only for the elements it shares with this one.

\textbf{History of the formulary.} The sacramentary witnesses are cited in H.~A.
Wilson's editions of the Gelasian (1894) and Gregorian (1915) sacramentaries,
read on their page images; the chant witnesses in gregorien.info's report of the
six manuscripts of Hesbert's \work{Antiphonale Missarum Sextuplex}, which was
not itself opened; and the lectionaries in the editions of Morin, Wilmart and
Migne. No manuscript, no critical edition of the sacramentaries and no Roman
Psalter was consulted, and none of these witnesses shows when this Gospel first
stood beside this Offertory.

\textbf{Chronology.} The Date cells in Scriptural Date and Location come from
the Scripture chronology corpus, through the record generated for this
formulary under its default profile. That record gives the three psalms only
the modern critical boundary for the Psalter, which it takes from the
introduction to the Psalms in the New American Bible, Revised Edition
(summarized, not quoted), and holds no traditional date for them. For the Gospel
and the Epistle it holds the traditional dates reported by the \work{Catholic
Encyclopedia}; for the Communion's psalm a superscription setting and a
prophetic referent; for the Offertory a relative date within the narrative;
and for the Introit two separate answers, one for each of its loci: the
\work{Catholic Encyclopedia}'s dates for the Book of Ecclesiasticus at the
antiphon's verse, and the Psalter's boundary at its psalm verse. The Gospel's
narrated event is undated.
The record and its sources are audited in \texttt{research/scope.md}.

\textbf{Rights.} The Latin of the formulary is published on the basis of its
public-domain antecedent, the Pustet Missal of 1862 (17 U.S.C. 103(b)); the
Douay--Rheims, the 1861 Cummiskey English, the Latin and Greek Fathers in Migne,
Zingerle's Hilary and the nineteenth-century English translations are in the
public domain, as are the 1909 English of the continuation of \work{The
Liturgical Year} and the 1927 English of Schuster. The Corpus Thomisticum text
of Aquinas on First Corinthians is a modern presentation and is quoted briefly. The New
American Bible introduction and the dated Ordos are cited, not reproduced.

\textbf{Records.} The appointed-text collation is in \texttt{propers/verified.md};
the search, loci, disagreements and negative results in
\texttt{research/scope.md}; and the reasoning behind the three readings in
\texttt{research/interpretations.md}. The research records were reviewed
independently before this study was written.
\par}

\section*{References}
% A starred heading sets no mark, so the Appendix that ends above would go on
% heading the pages of the bibliography. The heading stays starred, because
% check-content-preflight finds the References section by that exact form; the
% mark is set explicitly through the leaf-local \runninghead, which the web
% converter's macro whitelist carries.
\runninghead{References}
\addcontentsline{toc}{section}{References}
\label{sec:references}

{\footnotesize
\subsection*{Liturgical books and Bibles}
\begin{itemize}
\item \work{Missale Romanum ex decreto Sacrosancti Concilii Tridentini restitutum, Summorum Pontificum cura recognitum}, editio typica (Typis Polyglottis Vaticanis, 1962), CMAA facsimile, \url{https://media.churchmusicassociation.org/pdf/missale62.pdf}: \work{Rubricae generales} 111, 115, 127; \work{Rubricae generales Missalis} 427, 434, 494; pp.~409--411 (nos.~1668--1679).
\item \work{Missale Romanum}, editio iuxta typicam (New York: Benziger Brothers, 1962), pp.~404--405.
\item \work{Missale Romanum} (Ratisbon: Pustet, 1862), pp.~350--351, \latin{Dominica XVIII. post Pentecosten}.
\item \work{The Roman Missal, Translated into the English Language for the Use of the Laity} (Philadelphia: Eugene Cummiskey, 1861), pp.~444--445 (XVIII Sunday after Pentecost, Collect, Secret, Postcommunion).
\item The Holy Bible, Douay--Rheims version, revised by Bishop Richard Challoner (Project Gutenberg text): Ex 24:1--18; Ps 2:8; 95; 101:14--18; 121; 140:2; Ecclus 36:14--18; 50:29; Mt 8:28--9:9; 28:18; 1 Cor 1:1--13; 4:7; 16:8; 2 Par 6:21; 3 Kings 8:30.
\item \work{Biblia Sacra Vulgatae editionis} (Clementine Vulgate), eBible.org text: Ex 24; 29:41; Lev 8:11; Num 7:1; Ps 95; 101; 121; 140:2; Ecclus 36:18; Mt 9:1--8; Lk 2:32; 1 Cor 1:4--8.
\item \work{The Septuagint}, morphologically tagged Greek text (CATSS), at Ex 24:4--5, 17 and Sir 36:15--16; read for the comparison of wording and not reproduced.
\item H.~A. Wilson, ed., \work{The Gelasian Sacramentary: Liber Sacramentorum Romanae Ecclesiae} (Oxford, 1894), p.~232 (Liber III, sect.~XIV) and Liber I, sects.~LIX and XCIII.
\item H.~A. Wilson, ed., \work{The Gregorian Sacramentary under Charles the Great}, Henry Bradshaw Society 49 (1915), pp.~174--175 (sects.~CXXXV--CXXXVIII).
\item R.-J. Hesbert, \work{Antiphonale Missarum Sextuplex}, as reported by the gregorien.info chant database, \url{https://gregorien.info}: calendar day 823, \latin{Dominica XVIII post Pentecosten}, and the Gradual \latin{Timebunt gentes}.
\item Germain Morin, ``Le plus ancien \latin{Comes} ou lectionnaire de l'Église romaine'', \work{Revue Bénédictine} 27 (1910), pp.~62--63 (the Würzburg epistle list, item CLIII and note); ``Liturgie et basiliques de Rome au milieu du VII\textsuperscript{e} siècle'', \work{Revue Bénédictine} 28 (1911), pp.~315--316 (the Würzburg gospel list).
\item André Wilmart, ``Le \latin{Comes} de Murbach'', \work{Revue Bénédictine} 30 (1913), p.~50.
\item \work{Liber Comitis}, among the spuria of St Jerome, in Migne, PL 30, the Sundays of the seventh month.
\end{itemize}

\subsection*{Fathers, Doctors and saints}
\begin{itemize}
\item St Augustine, \work{Enarrationes in Psalmos}, trans. in Nicene and Post-Nicene Fathers, 1st ser., vol.~8 (CCEL text): on Ps~95, 1, 9--10; on Ps~101, s.~1, 16--18; on Ps~121, 2, 12--13. Latin in Migne, PL 37, cols.~1618--1619 (Ps~121, 2) and 1628 (Ps~121, 12).
\item St Augustine, \work{De consensu evangelistarum} II.25.57--58, trans. S.~D.~F. Salmond, Nicene and Post-Nicene Fathers, 1st ser., vol.~6 (CCEL text).
\item St Hilary of Poitiers, \work{Tractatus super Psalmos}, ed. A. Zingerle, CSEL 22 (Vienna, 1891): in Ps.~121, 1--5 and 13--14 (pp.~570--573, 578--579); \work{Commentarius in Matthaeum} VIII.4--8, Maurist text (Migne, PL 9, cols.~959--962).
\item St Ambrose, \work{Expositio Evangelii secundum Lucam} V.11--15 (Migne, PL 15, cols.~1357--1359; Latin Wikisource transcription).
\item St Bede, \work{In Lucae evangelium expositio} II, on Lk 5:21 and 5:24 (Migne, PL 92, cols.~388--389; web transcription).
\item St Jerome, \work{Commentarii in Matthaeum} I, on 9:1--8 (Migne, PL 26, cols.~54--56).
\item St John Chrysostom, \work{Homilies on Matthew} 29, Nicene and Post-Nicene Fathers, 1st ser., vol.~10 (CCEL text); \work{Expositio in Psalmum CXXI} (Migne, PG 55, cols.~347--351).
\item St John Chrysostom, \work{Homilies on First Corinthians} 2, and \work{Homilies on Hebrews} 16, Nicene and Post-Nicene Fathers, 1st ser., vols.~12 and 14, as delivered by New Advent (\url{https://www.newadvent.org/fathers/220102.htm}; \url{https://www.newadvent.org/fathers/240216.htm}).
\item Cassiodorus, \work{Expositio psalmorum}, in Ps.~95, 101 and 121, Latin text of Migne, PL 70 (Corpus Corporum and monumenta.ch delivery).
\item Theodoret of Cyrus, \work{Interpretatio in Psalmos}, in Ps.~95, 101 and 121 (Migne, PG 80, cols.~1647--1648, 1679--1680, 1879--1882); \work{Interpretatio in I ad Corinthios}, on 1:4--8 (Migne, PG 82, cols.~229--232).
\item \work{Expositiones in Psalmos} printed under the name of St Athanasius, on Ps~95 and 101 (Migne, PG 27, cols.~415--416, 429).
\item Theophylact of Ohrid, \work{In epistolam I ad Corinthios}, on 1:4--8 (Migne, PG 124, cols.~565--568).
\item Ambrosiaster, \work{Commentaria in Epistolam ad Corinthios Primam}, on 1:4--8 (Migne, PL 17; Latin Wikisource transcription).
\item St Thomas Aquinas, \work{Super Evangelium S.~Matthaei lectura} c.~9, in \work{Opera}, editio altera Veneta, vol.~3 (Venice, 1745), pp.~120--122.
\item St Thomas Aquinas, \work{Super I ad Corinthios} c.~1, lect.~1 (Corpus Thomisticum text).
\item St Robert Bellarmine, \work{A Commentary on the Book of Psalms}, trans. J. O'Sullivan (1866, abridged), read in the eCatholic2000 transcription of that translation: on Pss~95:8; 101:16; 121:1--2, 7.
\item Cornelius a Lapide, \work{Commentaria in Pentateuchum} (Antwerp, 1700), on Ex 24:4--8; \work{Commentaria in omnes D.~Pauli epistolas} (Antwerp, 1614), on 1 Cor 1:4--8.
\item St Anthony of Padua, sermon for the Nineteenth Sunday after Pentecost, in \work{Sermons for Sundays and Festivals}, trans. Paul Spilsbury; summarised only.
\end{itemize}

\subsection*{Liturgical commentators}
\begin{itemize}
\item Bl.\ Ildefonso Schuster, \work{The Sacramentary (Liber Sacramentorum)}, vol.~3 (London: Burns Oates \& Washbourne, 1927), pp.~167--170.
\item The continuation of Prosper Guéranger, \work{The Liturgical Year} (Dom Lucien Fromage), \work{Time after Pentecost}, vol.~2 (series vol.~11), trans. Laurence Shepherd, second edition (Stanbrook Abbey; London: Burns \& Oates, 1909), pp.~393--409.
\item Rupert of Deutz, \work{De divinis officiis} XII.18 (Migne, PL 170, col.~326).
\item William Durandus, \work{Rationale divinorum officiorum} VI.135 (Lyon, 1612), vol.~2.
\item Honorius Augustodunensis, \work{Gemma animae} IV.84--85 (Migne, PL 172, cols.~722--723).
\item Sicard of Cremona, \work{Mitrale} VIII.18 (Migne, PL 213, cols.~599--600).
\item Berno of Reichenau, \work{Libellus de quibusdam rebus ad Missae officium pertinentibus} V (Migne, PL 142, col.~1070).
\end{itemize}

\subsection*{Chronology and occurrence sources}
\begin{itemize}
\item \work{The Catholic Encyclopedia} (New York, 1907--1912): ``Biblical Chronology'' and ``General Chronology'' (vol.~3, 1908); ``Epistles to the Corinthians'' (vol.~4, 1908); ``Ecclesiasticus'' (vol.~5, 1909); ``Gospel of St.\ Matthew'' (vol.~10, 1911); ``St.\ Paul'' (vol.~11, 1911); ``The New Testament'' and ``Temple of Jerusalem'' (vol.~14, 1912).
\item United States Conference of Catholic Bishops, \work{New American Bible, Revised Edition}, introduction to the Psalms, on their dating; summarized, not reproduced.
\item Fraternité Saint-Pierre (France), \work{Ordo du mois}, entry for 27~September 2026; Institut du Christ Roi Souverain Prêtre (France), \work{Ordo}, entry for 27~September 2026.
\end{itemize}
\par}
